\section{The appointed texts}

The Latin follows the 1962 Missal (T), with its independently checked
Pustet 1862 antecedent (P). The primary English for the scriptural texts
is the Douay--Rheims, Challoner revision, from the registered Project
Gutenberg no. 1581 verse texts (B). Each scriptural appointment gives
those Bible verses first, unchanged. Lasance and Walsh's historical
liturgical form (L) follows separately, preserving the Missal's
introductions, adaptations and repetitions. The non-scriptural prayers,
Preface dialogue and Sanctus use L throughout. L is the revised-1945
expression in its later printing; short conclusions and the Introit's
doxology cue remain abbreviated as printed. Neither English source is
an approved English edition of the 1962 Missal. These inherited texts
retain their separate rights status. For edition details and textual
limits, see Scope and Qualifications.

\subsection{Introit -- Jer 29:11, 12, 14 adapted; Ps 84(85):2}
\label{proper-introit}
\properlocus{T p. 422, no. 1769; P p. 359; L p. 736.}

\begin{latinproper}
Dicit Dominus: Ego cogito cogitationes pacis, et non afflictionis:
invocabitis me, et ego exaudiam vos: et reducam captivitatem vestram de
cunctis locis. Ps. Benedixisti, Domine, terram tuam: avertisti captivitatem
Iacob. V. Gloria Patri.
\end{latinproper}

\begin{englishwitness}{Challoner: Jer 29:11}
For I know the thoughts that I think towards you, saith the Lord, thoughts
of peace, and not of affliction, to give you an end and patience.
\end{englishwitness}

\begin{englishwitness}{Challoner: Jer 29:12}
And you shall call upon me, and you shall go: and you shall pray to me,
and I will hear you.
\end{englishwitness}

\begin{englishwitness}{Challoner: Jer 29:14}
And I will be found by you, saith the Lord: and I will bring back your
captivity, and I will gather you out of all nations, and from all the
places to which I have driven you out, saith the Lord: and I will bring
you back from the place to which I caused you to be carried away captive.
\end{englishwitness}

\begin{englishwitness}{Challoner: Ps 84(85):2}
Lord, thou hast blessed thy land: thou hast turned away the captivity of Jacob.
\end{englishwitness}

The Introit shortens and rearranges the Jeremiah verses. Its complete
historical English, including the psalm verse and doxology cue, follows.

\begin{englishwitness}{Lasance/Walsh: historical liturgical form}
The Lord saith, I think thoughts of peace, and not of affliction; you
shall call upon Me and I will hear you; and I will bring back your
captivity from all places. Ps. Lord, Thou hast blessed Thy land: Thou
hast turned away the captivity of Jacob. V. Glory.
\end{englishwitness}

\subsection{Collect -- Excita, quaesumus}
\label{proper-collect}
\properlocus{T p. 423, no. 1770; P p. 359; L p. 737.}

\begin{latinproper}
Excita, quaesumus, Domine, tuorum fidelium voluntates: ut, divini operis
fructum propensius exsequentes; pietatis tuae remedia maiora percipiant.
Per Dominum.
\end{latinproper}

\begin{englishwitness}{Lasance/Walsh}
Arouse, we beseech Thee, O Lord, the wills of Thy faithful, that, by more
earnestly following after the fruit of the divine work, they may the more
abundantly partake of Thy mercies. Through our Lord.
\end{englishwitness}

\subsection{Epistle -- Col 1:9--14, with liturgical introduction}
\label{proper-epistle}
\properlocus{T p. 423, no. 1771; P p. 359; L p. 737.}

\begin{latinproper}
Lectio Epistolae beati Pauli Apostoli ad Colossenses.

Fratres: Non cessamus pro vobis orantes, et postulantes ut impleamini
agnitione voluntatis Dei, in omni sapientia et intellectu spiritali: ut
ambuletis digne Deo per omnia placentes: in omni opere bono fructificantes,
et crescentes in scientia Dei: in omni virtute confortati secundum potentiam
claritatis eius in omni patientia et longanimitate cum gaudio, gratias
agentes Deo Patri, qui dignos nos fecit in partem sortis sanctorum in
lumine: qui eripuit nos de potestate tenebrarum, et transtulit in regnum
Filii dilectionis suae, in quo habemus redemptionem per sanguinem eius,
remissionem peccatorum.
\end{latinproper}

\begin{englishwitness}{Challoner: Col 1:9--14}
Therefore we also, from the day that we heard it, cease not to pray for
you and to beg that you may be filled with the knowledge of his will,
in all wisdom and spiritual understanding:

That you may walk worthy of God, in all things pleasing; being fruitful
in every good work and increasing in the knowledge of God:

Strengthened with all might according to the power of his glory, in all
patience and longsuffering with joy,

Giving thanks to God the Father, who hath made us worthy to be partakers
of the lot of the saints in light:

Who hath delivered us from the power of darkness and hath translated us
into the kingdom of the Son of his love,

In whom we have redemption through his blood, the remission of sins:
\end{englishwitness}

The liturgical lesson replaces the letter's continuing opening with
``Brethren'' and names ``the will of God'' in place of ``his will.''

\begin{englishwitness}{Lasance/Walsh: historical liturgical form}
Lesson from the Epistle of blessed Paul the Apostle to the Colossians.

Brethren, We cease not to pray for you, and to beg that you may be filled
with the knowledge of the will of God, in all wisdom and spiritual
understanding; that you may walk worthy of God, in all things pleasing,
being fruitful in every good work, and increasing in the knowledge of God;
strengthened with all might according to the power of His glory, in all
patience and long-suffering with joy; giving thanks to God the Father, Who
hath made us worthy to be partakers of the lot of the saints in light;
Who hath delivered us from the power of darkness, and hath translated us
into the kingdom of the Son of His love, in Whom we have redemption through
His blood, the remission of sins.
\end{englishwitness}

\subsection{Gradual -- Ps 43(44):8--9}
\label{proper-gradual}
\properlocus{T p. 423, no. 1772; P p. 359; L pp. 737--738.}

\begin{latinproper}
Liberasti nos, Domine, ex affligentibus nos: et eos, qui nos oderunt,
confudisti. V. In Deo laudabimur tota die, et in nomine tuo confitebimur
in saecula.
\end{latinproper}

\begin{englishwitness}{Challoner: Ps 43(44):8--9}
But thou hast saved us from them that afflict us: and hast put them to
shame that hate us.

In God shall we glory all the day long: and in thy name we will give
praise for ever.
\end{englishwitness}

The chant begins \textit{Liberasti nos, Domine}, rendered by L as
``Thou hast delivered us, O Lord,'' where B has ``But thou hast saved us.''

\begin{englishwitness}{Lasance/Walsh: historical liturgical form}
Thou hast delivered us, O Lord, from them that afflict us: and hast put
them to shame that hate us. V. In God we will glory all the day: and in
Thy name we will give praise for ever.
\end{englishwitness}

\subsection{Alleluia -- Ps 129(130):1--2 adapted}
\label{proper-alleluia}
\properlocus{T p. 423, no. 1773; P p. 359; L p. 738.}

\begin{latinproper}
Alleluia, alleluia. V. De profundis clamavi ad te, Domine:
Domine, exaudi orationem meam. Alleluia.
\end{latinproper}

\begin{englishwitness}{Challoner: Ps 129(130):1--2}
Out of the depths I have cried to thee, O Lord:

Lord, hear my voice. Let thy ears be attentive to the voice of my supplication.
\end{englishwitness}

The chant asks God to hear the speaker's prayer in place of ``my voice''
and omits the rest of verse 2. L preserves this wording and the Alleluias.

\begin{englishwitness}{Lasance/Walsh: historical liturgical form}
Alleluia, alleluia. V. From the depths I have cried to Thee, O Lord:
Lord, hear my prayer. Alleluia.
\end{englishwitness}

\subsection{Gospel -- Mt 24:15--35, with liturgical introduction}
\label{proper-gospel}
\properlocus{T pp. 423--424, no. 1774; P pp. 359--360; L pp. 738--739.}

\begin{latinproper}
Sequentia sancti Evangelii secundum Matthaeum.

In illo tempore: Dixit Iesus discipulis suis: Cum videritis abominationem
desolationis, quae dicta est a Daniele propheta, stantem in loco sancto:
qui legit, intellegat: tunc qui in Iudaea sunt, fugiant ad montes: et qui
in tecto, non descendat tollere aliquid de domo sua: et qui in agro, non
revertatur tollere tunicam suam. Vae autem praegnantibus et nutrientibus
in illis diebus. Orate autem, ut non fiat fuga vestra in hieme vel sabbato.
Erit enim tunc tribulatio magna, qualis non fuit ab initio mundi usque modo,
neque fiet. Et nisi breviati fuissent dies illi, non fieret salva omnis
caro: sed propter electos breviabuntur dies illi. Tunc si quis vobis
dixerit: Ecce hic est Christus, aut illic: nolite credere. Surgent enim
pseudochristi et pseudoprophetae: et dabunt signa magna et prodigia, ita
ut in errorem inducantur (si fieri potest) etiam electi. Ecce praedixi
vobis. Si ergo dixerint vobis: Ecce in deserto est, nolite exire: ecce
in penetralibus, nolite credere. Sicut enim fulgur exit ab Oriente et
paret usque in Occidentem: ita erit et adventus Filii hominis.
Ubicumque fuerit corpus, illic congregabuntur et aquilae. Statim autem
post tribulationem dierum illorum sol obscurabitur, et luna non dabit
lumen suum, et stellae cadent de caelo, et virtutes caelorum commovebuntur:
et tunc parebit signum Filii hominis in caelo: et tunc plangent omnes
tribus terrae: et videbunt Filium hominis venientem in nubibus caeli
cum virtute multa et maiestate. Et mittet Angelos suos cum tuba et voce
magna: et congregabunt electos eius a quattuor ventis, a summis caelorum
usque ad terminos eorum. Ab arbore autem fici discite parabolam: cum
iam ramus eius tener fuerit et folia nata, scitis quia prope est aestas:
ita et vos cum videritis haec omnia, scitote quia prope est in ianuis.
Amen dico vobis, quia non praeteribit generatio haec, donec omnia haec
fiant. Caelum et terra transibunt, verba autem mea non praeteribunt.
\end{latinproper}

\begin{englishwitness}{Challoner: Mt 24:15--35}
When therefore you shall see the abomination of desolation, which was
spoken of by Daniel the prophet, standing in the holy place: he that
readeth let him understand.

Then they that are in Judea, let them flee to the mountains:

And he that is on the housetop, let him not come down to take any thing
out of his house:

And he that is in the field, let him not go back to take his coat.

And woe to them that are with child and that give suck in those days.

But pray that your flight be not in the winter or on the sabbath.

For there shall be then great tribulation, such as hath not been from
the beginning of the world until now, neither shall be.

And unless those days had been shortened, no flesh should be saved: but
for the sake of the elect those days shall be shortened.

Then if any man shall say to you, Lo here is Christ, or there: do not believe him.

For there shall arise false Christs and false prophets and shall shew
great signs and wonders, insomuch as to deceive (if possible) even the elect.

Behold I have told it to you, beforehand.

If therefore they shall say to you, Behold he is in the desert: go ye
not out. Behold he is in the closets: believe it not.

For as lightning cometh out of the east and appeareth even into the west:
so shall also the coming of the Son of man be.

Wheresoever the body shall be, there shall the eagles also be gathered together.

And immediately after the tribulation of those days, the sun shall be
darkened and the moon shall not give her light and the stars shall fall
from heaven and the powers of heaven shall be moved.

And then shall appear the sign of the Son of man in heaven. And then
shall all tribes of the earth mourn: and they shall see the Son of man
coming in the clouds of heaven with much power and majesty.

And he shall send his angels with a trumpet and a great voice: and they
shall gather together his elect from the four winds, from the farthest
parts of the heavens to the utmost bounds of them.

And from the fig tree learn a parable: When the branch thereof is now
tender and the leaves come forth, you know that summer is nigh.

So you also, when you shall see all these things, know ye that it is nigh,
even at the doors.

Amen I say to you that this generation shall not pass till all these things be done.

Heaven and earth shall pass: but my words shall not pass.
\end{englishwitness}

The Missal adds its liturgical introduction and omits ``therefore''
at the opening of verse 15. L supplies the complete historical lesson below.

\begin{englishwitness}{Lasance/Walsh: historical liturgical form}
Continuation of the holy Gospel according to St. Matthew.

At that time, Jesus said to His disciples, When you shall see the abomination
of desolation, which was spoken of by Daniel the prophet, standing in the
holy place; he that readeth, let him understand: then they that are in
Judea, let them flee to the mountains; and he that is on the house-top,
let him not come down to take anything out of his house; and he that is
in the field, let him not go back to take his coat. And woe to them that
are with child, and that give suck, in those days. But pray that your
flight be not in the winter, or on the sabbath: for there shall be then
great tribulation, such as has not been found from the beginning of the
world until now, neither shall be: and unless those days had been shortened,
no flesh should be saved; but for the sake of the elect, those days shall
be shortened. Then if any man shall say to you, Lo, here is Christ, or
there; do not believe Him: for there shall arise false Christs and false
prophets, and shall show great signs and wonders, insomuch as to deceive
(if possible) even the elect. Behold, I have told it to you beforehand:
if they therefore shall say to you, Behold, He is in the desert, go ye
not out; behold, He is in the closets, believe it not. For as lightning
cometh out of the east, and appeareth even into the west, so shall also
the coming of the Son of man be. Wheresoever the body shall be, there
shall the eagles also be gathered together. And immediately after the
tribulation of those days, the sun shall be darkened, and the moon shall
not give her light, and the stars shall fall from heaven, and the powers
of heaven shall be moved; and then shall appear the sign of the Son of
man in heaven, and then shall all the tribes of the earth mourn; and they
shall see the Son of man coming in the clouds of heaven with much power
and majesty: and He shall send His angels with a trumpet and a great
voice and they shall gather together His elect from the four winds, from
the farthest parts of the heavens to the utmost bounds of them. And from
the fig-tree learn a parable: when the branch thereof is now tender,
and the leaves come forth, you know that summer is nigh. So you also,
when you shall see all these things, know ye that it is nigh even at
the doors. Amen I say to you, that this generation shall not pass till
all these things be done. Heaven and earth shall pass away, but My words
shall not pass away.

\end{englishwitness}

The Creed follows the Gospel.

\subsection{Offertory -- Ps 129(130):1--2 adapted, repeated opening}
\label{proper-offertory}
\properlocus{T p. 424, no. 1775; P p. 360; L p. 739.}

\begin{latinproper}
De profundis clamavi ad te, Domine: Domine, exaudi orationem meam:
de profundis clamavi ad te, Domine.
\end{latinproper}

\begin{englishwitness}{Challoner: Ps 129(130):1--2}
Out of the depths I have cried to thee, O Lord:

Lord, hear my voice. Let thy ears be attentive to the voice of my supplication.
\end{englishwitness}

The Offertory has the Alleluia's shorter petition for God to hear the
speaker's prayer, then repeats the opening cry. It has no Alleluias.

\begin{englishwitness}{Lasance/Walsh: historical liturgical form}
From the depths I have cried out to Thee, O Lord: Lord, hear my prayer:
from the depths I have cried out to Thee, O Lord.
\end{englishwitness}

\subsection{Secret -- Propitius esto}
\label{proper-secret}
\properlocus{T p. 424, no. 1776; P p. 360; L p. 740.}

\begin{latinproper}
Propitius esto, Domine, supplicationibus nostris: et, populi tui
oblationibus precibusque susceptis, omnium nostrum ad te corda converte; ut,
a terrenis cupiditatibus liberati, ad caelestia desideria transeamus.
Per Dominum nostrum.
\end{latinproper}

\begin{englishwitness}{Lasance/Walsh}
Be propitious, O Lord, to our supplications, and accept the offerings and
prayers of Thy people: turn all our hearts unto Thee, that, being delivered
from earthly desires, we may pass on to the enjoyments of heaven.
Through our Lord.
\end{englishwitness}

\subsection{Preface of the Most Holy Trinity -- complete thanksgiving}
\label{proper-preface}
\properlocus{T p. 293, no. 1082 (appointed p. 424); P pp. 251--252; L pp. 773--774.}

\subsubsection{Dialogue}
\begin{latinproper}
V. Per omnia saecula saeculorum. R. Amen.

V. Dominus vobiscum. R. Et cum spiritu tuo.

V. Sursum corda. R. Habemus ad Dominum.

V. Gratias agamus Domino Deo nostro. R. Dignum et iustum est.
\end{latinproper}

\begin{englishwitness}{Lasance/Walsh}
P. World without end. R. Amen.

P. The Lord be with you. R. And with thy spirit.

P. Lift up your hearts. R. We have them lifted up unto the Lord.

P. Let us give thanks to the Lord our God. R. It is meet and just.
\end{englishwitness}

\begin{latinproper}
Vere dignum et iustum est, aequum et salutare, nos tibi semper et ubique
gratias agere: Domine sancte, Pater omnipotens, aeterne Deus: Qui cum
unigenito Filio tuo et Spiritu Sancto unus es Deus, unus es Dominus:
non in unius singularitate personae, sed in unius Trinitate substantiae.
Quod enim de tua gloria, revelante te, credimus, hoc de Filio tuo, hoc
de Spiritu Sancto sine differentia discretionis sentimus. Ut in
confessione verae sempiternaeque Deitatis, et in personis proprietas,
et in essentia unitas, et in maiestate adoretur aequalitas. Quam laudant
Angeli atque Archangeli, Cherubim quoque ac Seraphim: qui non cessant
clamare cotidie, una voce dicentes:
\end{latinproper}

\begin{englishwitness}{Lasance/Walsh}
It is truly meet and just, right and profitable, for us, at all times, and
in all places, to give thanks to Thee, O Lord, the holy One, the Father
almighty, the everlasting God: Who, together with Thine only-begotten Son
and the Holy Ghost, art one God, one Lord, not in the singleness of one
Person, but in the Trinity of one substance. For that which, according to
Thy revelation, we believe of Thy glory, the same we believe of Thy Son,
the same of the Holy Ghost, without difference or distinction; so that
in the confession of one true and eternal Godhead we adore distinctness
in persons, oneness in essence, and equality in majesty: Which the angels
praise, and the archangels, the cherubim also and the seraphim, who cease
not, day by day crying out with one voice, to repeat:
\end{englishwitness}

\subsubsection{Sanctus junction}
\begin{latinproper}
Sanctus, Sanctus, Sanctus Dominus Deus Sabaoth. Pleni sunt caeli et terra
gloria tua. Hosanna in excelsis. Benedictus qui venit in nomine Domini.
Hosanna in excelsis.
\end{latinproper}

\begin{englishwitness}{Lasance/Walsh}
Holy, holy, holy, Lord God of hosts. The heavens and the earth are full
of Thy glory. Hosanna in the highest. Blessed is He Who cometh in the
name of the Lord. Hosanna in the highest.
\end{englishwitness}

\subsection{Communion -- Mk 11:24 adapted}
\label{proper-communion}
\properlocus{T p. 424, no. 1777; P p. 360; L p. 740.}

\begin{latinproper}
Amen dico vobis, quidquid orantes petitis, credite quia accipietis,
et fiet vobis.
\end{latinproper}

\begin{englishwitness}{Challoner: Mk 11:24}
Therefore I say unto you, all things, whatsoever you ask when ye pray,
believe that you shall receive: and they shall come unto you.
\end{englishwitness}

L's Communion begins ``Amen, I say to you'' and ends ``it shall be done
unto you,'' where B begins ``Therefore I say unto you'' and ends
``they shall come unto you.''

\begin{englishwitness}{Lasance/Walsh: historical liturgical form}
Amen, I say to you, whatsoever you ask when you pray, believe that you
shall receive, and it shall be done unto you.
\end{englishwitness}

\subsection{Postcommunion -- Concede nobis}
\label{proper-postcommunion}
\properlocus{T p. 424, no. 1778; P p. 360; L p. 740.}

\begin{latinproper}
Concede nobis, quaesumus, Domine: ut per haec sacramenta quae sumpsimus,
quidquid in nostra mente vitiosum est, ipsorum medicationis dono curetur.
Per Dominum.
\end{latinproper}

\begin{englishwitness}{Lasance/Walsh}
Grant us, we beseech Thee, O Lord, that, through this sacrament which we
have received, whatever is evil in our hearts may be restored by its
gift of healing. Through our Lord.
\end{englishwitness}

The historical English's ``whatever is evil ... may be restored''
refers to the worshipper's cure: the Latin asks that the fault be healed.
The evil itself is not something to preserve.
